% ============================================================================
\section{Preliminaries}
\label{sec:setting}
% ============================================================================

\subsection{The MOPD objective}

There are $M$ tasks, each with a fixed specialist teacher. For a prompt $x$
from task $i$, the student samples one response $y$, and teacher $i$, denoted
by $\pi_{T_i}$, scores every position. At position $t$, the conditioning
context consists of the prompt $x$ and the student-generated response prefix
$y_{<t}$. If $y_t$ is sampled from the student's next-token distribution, the
usual reverse-KL token-gradient estimator is
\begin{multline*}
    \widehat h_i^{\mathrm{samp}}(x,y_{<t},y_t)\\[-0.4em]
    =\operatorname{sg}\!\left[
        \log \pi_\theta(y_t\mid x,y_{<t})
        -\log \pi_{T_i}(y_t\mid x,y_{<t})
      \right]\\[-0.2em]
    {}\times\nabla_\theta\log \pi_\theta(y_t\mid x,y_{<t}),
    \qquad y_t\sim\pi_\theta(\cdot\mid x,y_{<t}).
\end{multline*}
where $\operatorname{sg}$ stops gradients through the log-probability
difference. Thus the student is pulled toward the teacher on the tokens it
actually produced.

The teacher already computes scores over the vocabulary, and the student
already identifies its own high-probability tokens. Let
$\mathcal K_k(x,y_{<t})$ be the student's top-$k$ set at this context. For a
candidate token $a$, define
\[
    r_{i,a}(x,y_{<t})
    =\operatorname{sg}\!\left[
        \log \pi_\theta(a\mid x,y_{<t})
        -\log \pi_{T_i}(a\mid x,y_{<t})
      \right]
      \nabla_\theta\log \pi_\theta(a\mid x,y_{<t}).
\]
We compute the token gradient as
\begin{multline*}
    \widehat h_i^{\mathrm{top}k}(x,y_{<t},y_t)\\[-0.4em]
    =\sum_{a\in\mathcal K_k(x,y_{<t})}
       \pi_\theta(a\mid x,y_{<t})r_{i,a}(x,y_{<t})\\[-0.2em]
    {}+\mathbf{1}\!\left\{y_t\notin\mathcal K_k(x,y_{<t})\right\}
       r_{i,y_t}(x,y_{<t}),
    \qquad y_t\sim\pi_\theta(\cdot\mid x,y_{<t}).
\end{multline*}
At every fixed prefix this has the same expected gradient as the sampled-token
form; it integrates out which retained token was sampled, while the observed
token supplies an unbiased correction for the omitted tail (proof in
Appendix~\ref{app:opd_estimators}). It requires no additional teacher or
student pass because both distributions used by the expression are already
computed. In automatic differentiation, the top-$k$ membership and the
coefficients
$\pi_\theta(a\mid x,y_{<t})
 [\log\pi_\theta(a\mid x,y_{<t})-\log\pi_{T_i}(a\mid x,y_{<t})]$
are treated as stopped constants.

For each task, token contributions are averaged within a response and response
means are then averaged within a micro-batch; $L_i(\theta)$ denotes the
expectation of that task-level loss. Fixed positive weights, chosen before
training and summing to one, define the learning objective
\begin{equation}
    F(\theta)=\sum_{i=1}^M w_i L_i(\theta).
    \label{eq:fixed_multitask_objective}
\end{equation}
The fixed $w_i$ specify what the student learns, independently of how many
micro-batches estimate each term.

\subsection{Gradient of a mixed batch}

A micro-batch is a group of $b$ prompts from one task, processed in one
backward pass. An optimizer step accumulates $G$ micro-batches; task $i$
contributes the integer count $m_i$, the counts sum to $G$, and each count lies
between $m_{\min}$ and $m_{\max}$. The lower bound is at least two so that a
within-task sample variance exists on every step.

Let $g_{i,s}$ be the unweighted gradient of micro-batch $s$ from task $i$,
computed with the top-$k$ estimator above. Multiplying its loss by $w_i/m_i$
after the token and response averages gives the accumulated gradient
\begin{equation}
    A=\sum_{i=1}^M\frac{w_i}{m_i}\sum_{s=1}^{m_i}g_{i,s}.
    \label{eq:batch_gradient_observation}
\end{equation}
For fresh micro-batches that are conditionally independent across draws and
identically distributed within each task, with a common task mean and
covariance, let
$\mu_i(\theta)=\mathbb E[g_{i,s}\mid\theta]$ and
$\mu(\theta)=\sum_iw_i\mu_i(\theta)$. Every feasible allocation satisfies
$\mathbb E[A\mid m,\theta]=\mu(\theta)$, so changing the counts changes
precision but not the fixed weighted expected semi-gradient. When the rollout
distribution and all stopped coefficients are frozen at the current student,
$\mu$ is the gradient of the corresponding step-local surrogate for $F$. Any importance-ratio truncation is applied before the token
and response averages, and the factor $w_i/m_i$ is applied after them; in that
case the common expectation is for the same truncated surrogate.

A single mean over all valid tokens in the step would instead give a task more
influence when it contributes either more responses or longer responses. This
is the length-driven imbalance reported by Open-MOPD; response-level averaging
followed by $w_i/m_i$ weighting removes it by construction
\citep{gao2026openmopddiagnosingfixingcapability}. Table~\ref{tab:notation}
collects the notation.
